\begin{abstract}

Multi-teacher on-policy distillation (MOPD) trains a single student using specialized teachers, but the effects of loss normalization, vocabulary supervision, and optimizer dynamics on parameter updates and task performance are not well understood. We study these effects with a Qwen3-1.7B student and four teachers for mathematics, code, instruction following, and science. First, averaging losses over tokens rather than responses changes gradient directions even when total loss weights are equal between domains. Nevertheless, three normalization rules differ by only 0.14 percentage points in final mean score under student--teacher top-64 intersection supervision with Adam. Second, after 500 updates, training changes about 97\% of FP32 master weights but only 7.3--8.5\% of stored BF16 weights. Teacher-specific Adam updates appear aligned largely because they use the same accumulated moment estimates: subtracting the update that Adam would take with no new gradient reduces their mean cosine similarity from above 0.83 to below 0.01 on domain-specific inputs. Third, the intersection of the student's and teacher's top-64 token sets contains over 99.9\% of student probability mass on average at the evaluated checkpoints, and expanding supervision to the full vocabulary barely changes single-step update concentration. Yet we find no consistent performance advantage for top-64 over sampled-token supervision: at the joint-training endpoint, mathematics increases by 2.60 points while GPQA decreases by 1.77 points, with both paired 95\% confidence intervals including zero. Thus sparse BF16 changes do not imply sparse optimization, and aligned Adam updates do not imply aligned teacher gradients. Neither statistic can replace per-task evaluation when choosing a distillation objective.

\end{abstract}
